% Gerado por src/python/thesis/build_tab_fla_dinamica.py -- nao editar a mao.
% Regenerar com: make thesis-tables
% Fonte: results/tables/fla_model_comparison.csv
\begin{table}[!htb]
    \centering
    \caption{Capacidade preditiva dos descritores estáticos de paisagem sob validação deixa-uma-instância e deixa-uma-família.}
    \label{tab:fla_modelos}
    \resizebox{\textwidth}{!}{%
    \begin{tabular}{|l|l|c|c|c|c|}
        \hline
        \textbf{Modelo} & \textbf{Validação} & \textbf{$r_s$} & \textbf{$R^2$} & \textbf{BA} & \textbf{MCC} \\
        \hline
        Floresta aleatória (regressão) & Deixa-uma-instância & $0{,}374$ & $0{,}160$ & --- & --- \\
         & Deixa-uma-família & $0{,}380$ & --- & --- & --- \\
        \hline
        Floresta aleatória (classificação) & Deixa-uma-instância & --- & --- & $0{,}527$ & $0{,}059$ \\
         & Deixa-uma-família & --- & --- & $0{,}488$ & $-0{,}070$ \\
        \hline
        Regressão logística & Deixa-uma-instância & --- & --- & $0{,}728$ & $0{,}384$ \\
         & Deixa-uma-família & --- & --- & $0{,}439$ & $-0{,}163$ \\
        \hline
        SVC com núcleo RBF & Deixa-uma-instância & --- & --- & $0{,}690$ & $0{,}332$ \\
         & Deixa-uma-família & --- & --- & $0{,}488$ & $-0{,}070$ \\
        \hline
        $k$-vizinhos ($k=5$) & Deixa-uma-instância & --- & --- & $0{,}613$ & $0{,}280$ \\
         & Deixa-uma-família & --- & --- & $0{,}488$ & $-0{,}070$ \\
        \hline
    \end{tabular}
    }
    \par\smallskip
    \parbox{0.95\textwidth}{\footnotesize
    Modelos treinados sobre os 14 descritores estáticos de paisagem retidos, nas 51 instâncias rotuladas. $r_s$ é a correlação de Spearman entre ganho previsto e observado; BA é a acurácia balanceada e MCC o coeficiente de correlação de Matthews. A validação deixa-uma-família retira uma suíte inteira do treino, e é sob ela que a classificação binária deixa de generalizar: todos os modelos caem para a vizinhança do acaso, enquanto a regressão preserva o sinal. Os descritores estáticos servem, portanto, como prior sobre a magnitude do ganho, não como regra de decisão.
    }
\end{table}
